Machine-learning tools for finding security vulnerabilities in source code
typically represent each function as a graph of its syntax, control flow and
data dependencies, and train a neural network on that graph. Reported accuracy
is high and has improved steadily. It is also generally obtained by dividing
functions into training and test sets at random, which lets near-duplicates of
a test function, and often its own unfixed version, appear during training.

We rebuilt the pipeline over 131,903 real vulnerabilities from 949 open-source
projects, varying one factor at a time. Restricting only what the test set may
share with training lowers ROC-AUC from 0.796 to 0.585. Against the margin
above a coin flip, $(0.796-0.585)/(0.796-0.5)$, that is 71\% of the usable
signal: 47\% from duplicated functions, 24\% from shared projects, and 29\%
surviving every control.

That figure is not a property of the benchmark. Fixing the training recipe by
grid search on the validation partition alone, with test withheld until the
ranking was settled, raises the repository-disjoint score to 0.650 and lowers
the leakage share to 45\%. The recipe is worth $+0.032$ ($p=0.0001$), more than
any representational choice we measured, and it governs whether those choices
register at all: five-fold averaging is worth $+0.033$ under the default recipe
and $-0.002$ under the selected one, and bidirectional graph edges $+0.033$ and
$-0.006$. Across the twelve configurations searched, validation ROC-AUC predicts
test only weakly ($\rho=0.32$), so roughly a point of ROC-AUC is unreachable
without consulting the test partition.
